\section{Property Test Catalogue}
\label{app:properties}

This appendix restates the ten core invariants of Section~\ref{sec:invariants} as executable oracles. Each oracle is a precondition that constrains the generated input and a postcondition the operation under test must satisfy. The CDM enumeration universe (Section~\ref{sec:cdm}) supplies the generators. The framework draws random inputs, applies the operation, and checks the postcondition. A failure yields a minimal counterexample. The invariant prose is stated once in Section~\ref{sec:invariants}; this catalogue gives only the test forms.

Each row of the table reads the same way: given an input satisfying the precondition, the outcome must satisfy the postcondition.

\renewcommand{\arraystretch}{1.3}
\begin{center}
\begin{longtable}{@{}p{2.5cm} p{4.4cm} p{5.4cm}@{}}
\toprule
\textbf{Invariant} & \textbf{Precondition} & \textbf{Postcondition (oracle)} \\
\midrule
\endfirsthead
\toprule
\textbf{Invariant} & \textbf{Precondition} & \textbf{Postcondition (oracle)} \\
\midrule
\endhead

P1 Conservation & $\tau$ applied to state $s$, accepted & $\forall u:\ Q_{s'}(u)=\sum_w w(u)=0$ \\

P2 Atomic commitment & $\tau=\langle m_1,\dots,m_n\rangle$; move $m_k$ fails validation, $1\le k\le n$ & $\texttt{execute}(s,\tau)=s$; no partial state observable \\

P3 Referential integrity & $m$ a move in an accepted $\tau$ & $m.\mathrm{src}\in\mathcal{W}\ \wedge\ m.\mathrm{dst}\in\mathcal{W}\ \wedge\ m.\mathrm{unit}\in\mathcal{U}$ \\

P4 Log monotonicity & log $L=\langle e_1,\dots,e_n\rangle$ after any operation sequence & $\forall i>1:\ e_i.\mathrm{prevHash}=H(e_{i-1})$; no $e_j$, $j<n$, rewritten (append-only, hash-chained) \\

P5 Transaction idempotency & state $s$, transaction $\tau$ & $\texttt{execute}(\texttt{execute}(s,\tau),\tau)=\texttt{execute}(s,\tau)$ (dedup by \texttt{tx\_id}) \\

P6 Lifecycle idempotency & $f(u,\mathit{st},d)=(\mathit{moves},\mathit{st}')$; $\mathit{st}'$ already reflects the event & $f(u,\mathit{st}',d)=(\varnothing,\mathit{st}')$; second call emits no moves \\

P7 Ledger isolation & $m$ a move in virtual ledger $\mathcal{L}_v$ & $m.\mathrm{src}\in\mathcal{L}_v\ \wedge\ m.\mathrm{dst}\in\mathcal{L}_v$; no move crosses the virtual/real boundary \\

P8 Snapshot consistency & time $t$ in the log; $s_t$ the state at $t$ & $\forall w,u:\ \texttt{clone\_at}(t).\texttt{balance}(w,u)=s_t.\texttt{balance}(w,u)\ \wedge\ \texttt{clone\_at}(t).\texttt{unit\_state}(u)=s_t.\texttt{unit\_state}(u)$ \\

P9 Lifecycle purity & lifecycle function $f$, fixed inputs $(u,\mathit{st},d)$ & repeated evaluation yields identical $(\mathit{moves},\mathit{st}')$; no effect beyond the return \\

P10 PnL path-independence & price vectors $P_{t_0},P_{t_1}$; any transaction sequence on $[t_0,t_1]$ & $\mathrm{PnL}_{\text{price}}+\mathrm{PnL}_{\text{flow}}=V_{t_1}-V_{t_0}$, independent of the intermediate path \\

\bottomrule
\end{longtable}
\end{center}
\renewcommand{\arraystretch}{1.0}

\begin{remark}[P8 and the projection discipline]
The P8 oracle reconstructs \emph{both} balances and unit state from the log. \texttt{unit\_state} is read from UnitStatus, a materialised projection of the immutable event log---a read cache the log always rebuilds, never a separate source of truth (Section~\ref{sec:lifecycle}). The oracle is the executable form of that discipline: a catamorphism---replaying the events up to $t$ rebuilds the exact value that held then---so $\texttt{clone\_at}(t).\texttt{unit\_state}(u)=s_t.\texttt{unit\_state}(u)$ holds by construction, and a failure exposes any drift between the cache and the log. The observation \texttt{sameLedger} makes through \texttt{unitStatus} covers \texttt{usBasis}, so P8's generator universe must include boundary transactions with retro-effective $t_{\mathrm{eff}}$.
\end{remark}

\noindent\textbf{Totality oracles.} Two further oracles are global. The same generator exercises them, but they are not numbered among the core ten:
\begin{itemize}
\item \emph{Totality.} For every generated input $i$, $\texttt{execute}(s,i)$ returns success or a well-typed error; it never raises an unhandled exception.
\item \emph{Valid transitions only.} For every $(\text{product type},\text{state},\text{event})$ triple absent from the product's transition table, the lifecycle function returns an error. This is the executable witness that illegal lifecycle transitions are unrepresentable (Section~\ref{sec:lifecycle}).
\end{itemize}

\noindent\textbf{PnL inputs.} $V_t$ and the price vectors $P_t(u)$ that drive P10 are produced by the valuation projection (Section~\ref{sec:valuation}). The state-aware price function is specified in Appendix~\ref{app:pricing-coordination}.

\noindent\textbf{SBL, obligation, basis, producer, and vocabulary extensions.} Section~\ref{sec:sbl-invariants} adds ten invariants P11--P20 specific to Securities Borrowing and Lending (collateralisation, lender ownership invariance, partial-return and settlement monotonicity). Section~\ref{sec:obligation-invariants} adds P21--P23 for obligation liveness. Section~\ref{sec:state-basis} adds P24, basis tip agreement, and P25, ingest-door soundness. Section~\ref{sec:sbl-locates} adds P26, locate-capacity admission. Section~\ref{sec:invariants-producer} adds five more: P27, producer agreement (oracle \texttt{recomputeOK}, with the terms-leg oracle \texttt{propTermsRecompute}); P28, cause committed (\texttt{pCauseCommitted}); P29, schedule totality, the witness of \hyperref[cond:C14]{C14} (\texttt{propScheduleTotalGate}, quantifying over the closed \texttt{CAClass}); P30, dimension invariance (\texttt{propDimInvariance}, with the composite-kind oracle \texttt{dkDivSplitOK}); and P31, kind totality at the ingest door (\texttt{pKindTotal}, the P25 shape).

These extensions never relax a core invariant. Every transaction satisfying P11--P23 also satisfies P1--P10. P24 quantifies over log prefixes rather than single transactions, P26 over locate-confirmation commit positions, and P27 over committed logs with a recomputing implementation; each narrows the accepted logs without weakening any oracle.

The executable forms of P11--P23 follow the same precondition/postcondition shape. Those of P24--P31 are given separately: P24 below, P25 with the market-data contract (\S\ref{sec:basis-contract}), P26 as the log-replay check of Section~\ref{sec:sbl-locates}, and P27--P31 as the named suite oracles of Section~\ref{sec:invariants-producer} (\texttt{reference/Ledger.hs}, Part~L and the gates they witness). Three conformance-tier properties --- P32, confirmation gate; P33, elective aggregate invariance; P34, basket joint consumption --- are specified in Section~\ref{sec:invariants-producer}; their executable forms live where \texttt{transport} lives, outside the reference.

P24 asserts that two computations of a unit's basis tip agree: the fold that applies boundary events in booking order, and the projection that ranks them in effective order. The oracle checks the agreement at every accepted prefix of the log.

\renewcommand{\arraystretch}{1.3}
\begin{center}
\begin{longtable}{@{}p{2.5cm} p{4.4cm} p{5.4cm}@{}}
\toprule
\textbf{Invariant} & \textbf{Precondition} & \textbf{Postcondition (oracle)} \\
\midrule
\endfirsthead
\toprule
\textbf{Invariant} & \textbf{Precondition} & \textbf{Postcondition (oracle)} \\
\midrule
\endhead

P24 Basis tip agreement & a log $L$; any prefix $L_k=\langle\tau_1,\dots,\tau_k\rangle$ whose replay from the empty ledger is accepted & $\forall u$ registered in $L_k$: $\texttt{usBasis}(u)=\mathrm{effTip}(u,L_k)$---the maximum, in the lexicographic order $(t_{\mathrm{eff}},\mathit{prec},\mathit{bid})$, of the boundary events $L_k$ carries on $u$, or $\mathrm{Origin}(u)$ if it carries none. The booking-order status fold and the effective-order chain projection agree at every prefix \\

\bottomrule
\end{longtable}
\end{center}
\renewcommand{\arraystretch}{1.0}

\begin{remark}[P24 and the time-travel claim]
The basis-time-travel claim of \S\ref{sec:basis-time-travel} is discharged by P8 together with P24: each accepted-prefix check is a $\texttt{clone\_at}(k)$, and \texttt{sameLedger} observes \texttt{usBasis} through \texttt{unitStatus}. No-regression is a corollary, not a separate test: the maximum of a growing committed set is monotone. The typed oracle is \texttt{p24} (\texttt{reference/Ledger.hs}, Part~L); its recomputation \texttt{effTip} reads the transactions' own boundary events, never the \texttt{ledgerBounds} cache.
\end{remark}

\medskip
\noindent\textbf{Typed oracles.} Each oracle above has a typed counterpart. Its precondition and postcondition are both total, so the generator and the checker share one type. The operation under test returns a typed \texttt{Outcome}. \texttt{Accepted} is the \emph{only} constructor that carries a ledger, so a partial (``half-committed'') ledger is not representable. That absence is the \emph{Totality} oracle and the ``no partial state'' half of P2, made structural rather than tested.

The listing fixes the shared vocabulary: the outcome type, the property shape, and the observational equality the oracles use.

\begin{lstlisting}[language=Haskell]
-- execute = the canonical (Sec. 4) three-home transaction-apply, lifted to
-- return Outcome; supplied by the implementation, not a symbol defined here.
data Outcome = Accepted Ledger | Rejected TxError   -- no third "crashed" case

-- A property: a precondition constraining the generated input and a
-- postcondition over the outcome of running the operation on it. Both total.
data Property i o = Property
  { pre  :: Ledger -> i -> Bool
  , post :: Ledger -> i -> o -> Bool }

data TxInput = TxInput { tiTx :: OracleTx, tiScope :: Scope }
data Scope   = Scope   { scWallets :: [WalletId], scUnits :: [UnitId] }

registered :: Ledger -> UnitId -> Bool
registered l u = maybe False (const True) (productTerms l u)

-- Sum a wallet's holding of u across a wallet set. Summation is a monoid
-- homomorphism into the Qty group -- the empty set sums to mempty (=0) with no
-- special case, so a zero-holder unit conserves vacuously (no divide-by-holders).
unitTotal :: Ledger -> [WalletId] -> UnitId -> Qty
unitTotal l ws u = foldMap (\w -> maybe mempty psBalance (position l w u))
                           (Set.toList (Set.fromList ws))

-- Equality of ledgers BY OBSERVATION over the scope -- the only equality the
-- oracles need; the abstract Ledger exports no structural Eq.
sameLedger :: Scope -> Ledger -> Ledger -> Bool
sameLedger sc a b =
     and [ position a w u == position b w u | w <- scWallets sc, u <- scUnits sc ]
  && and [ unitStatus a u == unitStatus b u | u <- scUnits sc ]
\end{lstlisting}

\noindent Five oracles test the same operation, \texttt{execute}, and share the one \texttt{Property TxInput Outcome} shape. Each postcondition that asserts a positive fact is guarded on \texttt{Accepted}, so it claims nothing about a rejected transaction.

\begin{lstlisting}[language=Haskell]
-- P1 Conservation. On any accepted transaction every unit's holdings sum to
-- zero across the scope (closed system); a rejected one asserts nothing.
p1 :: Property TxInput Outcome
p1 = Property (\_ _ -> True) $ \_ (TxInput _ sc) o -> case o of
  Rejected _ -> True
  Accepted l -> all (\u -> unitTotal l (scWallets sc) u == mempty) (scUnits sc)

-- P2 Atomic commitment. Any move failing validation rejects the whole txn;
-- no partial ledger exists (Outcome cannot carry one).
p2 :: Property TxInput Outcome
p2 = Property (\l (TxInput tx _) -> any (invalidMove l) (otMoves tx))
              (\_ _ o -> isRejected o)

-- P3 Referential integrity. Every move in an accepted txn names a registered
-- unit and in-scope wallets.
p3 :: Property TxInput Outcome
p3 = Property (\_ _ -> True) $ \_ (TxInput tx sc) o -> case o of
  Rejected _ -> True
  Accepted l -> all (\m -> registered l (mUnit m)
                        && inScope sc (mFrom m) && inScope sc (mTo m)) (otMoves tx)

-- P7 Ledger isolation. With Scope = the virtual ledger's wallets, no move
-- crosses its boundary -- a property of the moves, independent of the outcome.
p7 :: Property TxInput Outcome
p7 = Property (\_ (TxInput tx _) -> not (null (otMoves tx)))
     $ \_ (TxInput tx sc) _ ->
         all (\m -> inScope sc (mFrom m) && inScope sc (mTo m)) (otMoves tx)

-- P5 Transaction idempotency (dedup by tx_id). Re-running an accepted txn is a
-- no-op; this oracle re-runs the operation, so it takes it directly.
idempotentTx :: (Ledger -> OracleTx -> Outcome)
             -> Scope -> Ledger -> OracleTx -> Bool
idempotentTx exec sc l tx = case exec l tx of
  Rejected _  -> True
  Accepted l' -> case exec l' tx of Accepted l'' -> sameLedger sc l' l''
                                    Rejected _   -> False
\end{lstlisting}

\noindent The remaining five test four other operations: the log append, the snapshot rebuild, the pure lifecycle handler, the PnL decomposition. Forcing them into the \texttt{Property} shape would buy nothing. Each is therefore a bespoke total predicate (the restraint rule: stop at the cheaper guarantee).

\begin{lstlisting}[language=Haskell]
-- P4 Log monotonicity. The log is append-only and hash-chained BY CONSTRUCTION
-- (built only by `append`, which stamps the predecessor's hash; no rewrite or
-- delete is exported). The oracle re-checks the link; a break needs API-tamper.
data Hashed = Hashed { hPrev :: Hash, hPayload :: Transaction, hSelf :: Hash }
chainOK :: [Hashed] -> Bool
chainOK es = and (zipWith (\a b -> hPrev b == hSelf a) es (drop 1 es))

-- P8 Snapshot consistency. cloneAt t rebuilds the state at t by replaying the
-- events up to t (a catamorphism -- replaying the log always rebuilds the exact
-- value that held then), so snapshot must equal s_t by observation (both
-- balances and unit state -- sameLedger covers both).
p8 :: Scope -> Ledger -> Ledger -> Bool          -- cloneAt t  vs  s_t
p8 = sameLedger

-- A lifecycle handler is a PURE TOTAL function. Purity (no IO in the type) IS
-- P9: re-evaluation cannot differ and there is no effect to leak.
type LifecycleFn st ev = UnitId -> st -> ev -> ([Move], st)

-- P6 Lifecycle idempotency. Feeding the handler its own output state emits no
-- further moves and fixes the state.
idempotentL :: (Eq st) => LifecycleFn st ev -> UnitId -> st -> ev -> Bool
idempotentL f u st ev = let (_, st') = f u st ev in f u st' ev == ([], st')

-- Valid-transitions-only. For every (product, state, event) absent from the
-- transition table, the guarded handler returns an error: not in table => Left.
-- An illegal transition is thus not representable as success.
validTransitionsOnly
  :: (p -> st -> ev -> Bool)                         -- in the table?
  -> (p -> st -> ev -> Either TxError ([Move], st))  -- the guarded handler
  -> p -> st -> ev -> Bool
validTransitionsOnly inTable h p st ev
  | inTable p st ev = True
  | otherwise       = case h p st ev of Left _ -> True; Right _ -> False

-- P10 PnL path-independence. V_t and the price vectors P_t(u) come from the
-- valuation projection; the decomposition closes regardless of the
-- intermediate transaction path.
p10 :: Cash -> Cash -> Cash -> Cash -> Bool      -- pnlPrice pnlFlow V_t0 V_t1
p10 pnlPrice pnlFlow v0 v1 = pnlPrice <> pnlFlow == v1 <> cashNeg v0
\end{lstlisting}

\noindent P24 is log-prefix-shaped, like P4 and P8. The check replays every prefix from the empty ledger and compares the cached tip with a recomputation from the transactions' own boundary events. Its oracle and the spec function it recomputes are verbatim from the reference (Part~L):

\begin{lstlisting}[language=Haskell]
-- | The effective-order tip a log prefix mandates for u, recomputed from the
--   transactions' own BoundaryEvents -- INDEPENDENT of the ledger's ledgerBounds
--   cache, so a drifted cache cannot vouch for itself. The spec function of the
--   basis-chain definition (sec:state-basis), written once, here: the maximum,
--   in the lexicographic order (t_eff, prec, bid), of the boundary events the
--   prefix carries on u, or Origin u if it carries none.
effTip :: UnitId -> [Transaction] -> BasisId
effTip u txs =
  case Set.lookupMax (Set.fromList
         [ (bevTEff be, bevPrec be, bevId be)
         | tx <- txs, txUnit tx == u, Just be <- [txBoundary tx] ]) of
    Nothing          -> Origin u
    Just (_, _, bid) -> Boundary bid

-- P24 Basis tip agreement: at EVERY accepted prefix of the log, the booking-order
-- status fold (replay) and the effective-order chain projection (effTip) agree on
-- usBasis for every registered unit. Quantifying over prefixes makes each check a
-- clone_at(k), so the retro-effective time-travel claim (sec:basis-time-travel)
-- rides on this oracle plus the existing P8 (sameLedger already observes usBasis
-- through unitStatus, whose structural Eq covers all four fields). No-regression
-- is a corollary, not a test: the max of a growing committed set never recedes.
-- Wrongful ADMISSION of a non-tip SetBasis fails this oracle at that prefix;
-- wrongful REJECTION is the vacuous side, as in P1/P2's split -- consistent with
-- catalogue precedent. Log-prefix-shaped like P4/P8, not Property-shaped: the
-- claim quantifies over prefixes of a log, not over single transactions.
-- (Generators: BoundaryId is abstract by design, so boundary transactions for
-- this oracle are compiled IN-MODULE, minting deterministic content addresses --
-- no production caller gains a minting door.)
p24 :: [Transaction] -> Bool
p24 txs = all prefixOK [ take k txs | k <- [0 .. length txs] ]
  where
    prefixOK pfx = case replay pfx emptyLedger of
      Left _  -> True   -- a rejected prefix asserts nothing (rejection is P2's half)
      Right l -> all (agrees l pfx) (Set.toList (Set.fromList (map txUnit pfx)))
    agrees l pfx u = fmap usBasis (unitStatus l u) == Just (effTip u pfx)
\end{lstlisting}

\noindent The full self-contained listing is collected in \texttt{reference/Ledger.hs}. It supplies the helpers \texttt{inScope}, \texttt{isRejected}, \texttt{invalidMove}, the \texttt{Move}/\texttt{OracleTx} types, the \texttt{TxError} and \texttt{Hash} types, the \texttt{Cash} group inverse \texttt{cashNeg}, and the \texttt{total} coverage witness. It reuses the \texttt{Ledger} library and its total accessors.

